% ------------------------------------------------------------------------------------
%                              ESERCIZI SISTEMI LTI
% ------------------------------------------------------------------------------------

\chapter{Esercizi in aula - Sistemi}

\section{Esercizio 1}

Consideriamo il sistema mostrato in figura~\ref{fig:Es_01_01_Esercizio1_Sistema}.

\begin{figure}[!ht]
    \centering
    \begin{tikzpicture}[
        >=Stealth, % Punta delle frecce
        thick,     % Spessore delle linee
        node distance=2cm
        ]
        % --- Definizione dei nodi principali ---

        % Nodo di somma (il cerchio con + e -)
        \node[circle, draw, inner sep=1pt, minimum size=7mm] (sum) at (4.5,0) {};
        \node at ([shift={(-0.15, 0.12)}]sum.center) {\(+\)};
        \node at ([shift={(0, -0.15)}]sum.center) {\(-\)};

        % Blocco Integratore
        \node[draw, rectangle, minimum width=1.8cm, minimum height=1.2cm] (int) at (7.5,0) {\huge \(\int\)};

        % Blocco Ritardo (RIT)
        \node[draw, rectangle, minimum width=1.2cm, minimum height=0.8cm] (delay) at (2.5,-1.8) {RIT};
        % Etichetta sotto il blocco ritardo
        \node[below=0.1cm of delay] {\(T_0=1\)};

        % --- Tracciamento delle linee e frecce ---

        % Linea di ingresso x(t)
        \draw[->] (0,0) node[above right] {\(x(t)\)} -- (sum);

        % Nodo di derivazione (il pallino nero)
        \fill (2.5,0) circle (1.5pt);

        % Freccia dal sommatore all'integratore
        \draw[->] (sum) -- node[above] {\(z(t)\)} (int);

        % Freccia dall'integratore all'uscita
        \draw[->] (int) -- node[above] {\(y(t)\)} (10,0);

        % --- Percorso di retroazione (feedback) ---

        % Freccia che scende dal nodo di derivazione al blocco ritardo
        \draw[->] (2.5,0) -- (delay.north);

        % Freccia che esce dal ritardo, va a destra e poi sale verso il sommatore
        \draw[->] (delay.east) -| (sum.south);
    \end{tikzpicture}
    \caption{Sistema d'esempio}
    \label{fig:Es_01_01_Esercizio1_Sistema}
\end{figure}

\begin{exe}
    Vogliamo studiare il sistema per capire:
    \begin{itemize}
        \item Se è \textbf{LTI}.
        \item Qual è la sua \textbf{risposta impulsiva} e la sua
            \textbf{ripsosta in frequenza}.
    \end{itemize}
\end{exe}

Cominciamo l'analisi considerando solamente il \textbf{primo blocco}:
\[
    \begin{aligned}
        z(t) &= x(t) - x(t-T_{0}) \\ 
             &= x(t) - x(t-1)
    \end{aligned}
\]

Verifichiamo che il sistema sia \textbf{lineare}, considerando come
ingresso $x(t) = \alpha x_{1}(t) + \beta x_{2}(t)$
\[
    \begin{aligned}
        \alpha x_{1}(t) &\longrightarrow \alpha[ x_{1}(t) - x_{1}(t-1) ] = y_{1}(t) \\ 
        \beta x_{2}(t) &\longrightarrow \beta[x_{2}(t) - x_{2}(t-1)] = y_{2}(t) \\
        x(t) &\longrightarrow [\alpha x_{1}(t) + \beta x_{2}(t)] - [\alpha x_{1}(t-1) + \beta x_{2}(t-1)] \\ 
             &= y_{1}(t) + y_{2}(t)
    \end{aligned}
\]

Ora verifichiamo anche che il sistema sia \textbf{tempo invariante}:
\[
    \begin{aligned}
        x(t) &\longrightarrow z(t) = x(t) - x(t-1) \\ 
        x(t-t_{0}) &\longrightarrow w(t) = x(t-t_{0}) - x(t-1-t_{0}) \\ 
                   &= x(t-t_{0}) - x(t-t_{0}-1) = z(t-t_{0})
    \end{aligned}
\]
\begin{remark}
    Per cui il primo blocco è un \textbf{sistema LTI}, e quindi possiamo 
    rappresentarlo come un \textbf{filtro}, come mostrato in 
    figura~\ref{fig:Es_01_02_Esercizio1_SistemaFiltro}.
\end{remark}

\begin{figure}[!ht]
    \centering
    \begin{tikzpicture}[
        >=Stealth, % Punta delle frecce
        thick,     % Spessore delle linee
        block/.style={draw, rectangle, minimum width=2cm, minimum height=1.2cm}
        ]
        % --- Nodi principali ---

        % Blocco Filtro con risposta h_1(t)
        \node[block] (filt) at (3,0) {\(h_1(t)\)};

        % Blocco Integratore
        \node[block] (int) at (7.5,0) {\huge \(\int\)};

        % --- Tracciamento delle linee e frecce ---

        % Linea di ingresso x(t)
        \draw[->] (0,0) -- (filt) node[midway, above] {\(x(t)\)};

        % Freccia dal filtro all'integratore
        \draw[->] (filt) -- (int) node[midway, above] {\(z(t)\)};

        % Freccia dall'integratore all'uscita
        \draw[->] (int) -- (10.5,0) node[midway, above] {\(y(t)\)};
    \end{tikzpicture}
    \caption{Sostituzione del primo blocco con un filtro}
    \label{fig:Es_01_02_Esercizio1_SistemaFiltro}
\end{figure}

Calcoliamo ora la \textbf{risposta impulsiva} $h_{1}(t)$ di questo primo blocco.
Per farlo, basta considerare l'\textbf{uscita del sistema} 
quando l'\textbf{ingresso} è l'\textbf{impulso di Dirac}, e risulta
banalmente:
\[
    h_{1}(t) = \delta(t) - \delta(t-1)
\]
Per cui la \textbf{risposta in frequenza} è:
\[
    H_{1}(f) = 1 - e^{-j 2\pi f}
\]
Notiamo che:
\[
    H_{1}(-f) = 1 - e^{j 2\pi f} = \overline{H_{1}(f)} 
\]

\begin{remark}
    La risposta in frequenza ha \textbf{simmetria hermitiana}, esattamente
    come ci aspettavamo visto che $h_{1}(t)$ è \textbf{reale}.
\end{remark}

Passiamo ora al \textbf{secondo blocco} del sistema:
\[
    y(t) = \int_{-\infty}^t z(\tau ) d \tau
\]

Come prima, vogliamo verificare se il sistema è \textbf{LTI}. 
Per prima cosa notiamo che, essendo l'integrale un 
\textit{operatore lineare}, il sistema è \textit{sicuramente lineare}. 
Verifichiamo allora se è anche \textit{tempo invariante}:
\[
    \begin{aligned}
        z(t) &\longrightarrow y(t) = \int_{-\infty}^t z(\tau ) d \tau \\ 
        z(t-t_{0}) &\longrightarrow \int_{-\infty}^{t} z(\tau -t_{0}) d \tau \\
                   & u \coloneqq \tau - t_{0} \text{ , } du = d\tau 
                   \text{ , } \tau \in (-\infty, t] \implies u \in (-\infty, t-t_{0}] \\ 
                   &= \int_{-\infty}^{t-t_{0}} z(u) du \\ 
                   &= y(t-t_{0})
    \end{aligned}
\]
Appurato che anche il secondo blocco è un sistema \textbf{LTI}, possiamo
rappresentare il \textit{sistema complessivo} come una \textbf{cascata di due sistemi LTI},
come mostrato in figura~\ref{fig:Es_01_03_Esercizio1_RappresentazioneAlternativaSistema}.

\begin{figure}[!ht]
    \centering
    \begin{tikzpicture}[
        >=Stealth, % Punta delle frecce
        thick,     % Spessore delle linee
        block/.style={draw, rectangle, minimum width=2cm, minimum height=1.2cm}
        ]
        % --- Nodi principali ---

        % Primo blocco (Filtro h1)
        \node[block] (filt1) at (2.5,0) {\(h_1(t)\)};

        % Secondo blocco (Filtro h2)
        \node[block] (filt2) at (7,0) {\(h_2(t)\)};

        % --- Tracciamento delle linee e frecce ---

        % Linea di ingresso x(t) al primo blocco
        \draw[->] (0,0) -- (filt1) node[midway, above] {\(x(t)\)};

        % Freccia tra i due blocchi (segnale intermedio)
        \draw[->] (filt1) -- (filt2) node[midway, above] {\(z(t)\)};

        % Freccia in uscita dal secondo blocco
        \draw[->] (filt2) -- (9.5,0) node[midway, above] {\(y(t)\)};
    \end{tikzpicture}
    \caption{Rappresentazione alternativa del sistema}
    \label{fig:Es_01_03_Esercizio1_RappresentazioneAlternativaSistema}
\end{figure}

A questo punto, come prima, vogliamo calcolare la \textbf{risposta impulsiva}
del blocco:
\[
    h_{2}(t) = \int_{-\infty}^t \delta(\tau ) d \tau 
\]
Notiamo che, per definizione dell'impulso, l'integrale è pari a \textit{uno} 
se $t\ge 0$ e \textit{zero altrimenti}.
Per cui abbiamo:
\[
    h_{2}(t) = \begin{cases}
        1 & t\ge 0 \\ 
        0 & t < 0
    \end{cases} = \mathbf{1}(t)
\]


Allora che la \textbf{risposta complessiva} del sistema è:
\[
    \begin{aligned}
        y(t) &= h_{2}(t) \ast z(t) \\ 
             &= h_{2}(t) \ast [ h_{1}(t) \ast x(t)] \\ 
             &= [h_{1}(t) \ast h_{2}(t)] \ast x(t)
    \end{aligned}
\]

Per cui possiamo immaginare il sistema come un \textbf{filtro equivalente}
con \textbf{risposta impulsiva}:
\[
    \begin{aligned}
        h(t) &= h_{1} \ast h_{2}(t) \\ 
             &= [\delta(t) - \delta(t-1)] \ast \mathbf{1}(t) \\ 
             &= \mathbf{1}(t) - \mathbf{1}(t-1)
    \end{aligned}
\]

Vediamo $h(t)$ in figura~\ref{fig:Es_01_04_Esercizio1_RispostaImpulsivaEquivalente}
per avere un'intuizione su come riscriverla.

\begin{figure}[!ht]
    \centering
    \begin{tikzpicture}
        \begin{axis}[
            width=0.75\textwidth,
            height=5cm,
            axis lines=middle,
            xlabel={$t$},
            ylabel={$h(t)$},
            xmin=-1.5, xmax=2.5,
            ymin=-0.2, ymax=1.2,
            xtick={-1,0,1,2},
            ytick={0,1},
            grid=major,
            grid style={dashed, gray!30},
            xlabel style={below right},
            ylabel style={above left},
            ]

            % Linee tratteggiate per evidenziare le discontinuità (salti)
            \draw[dashed, gray!60] (axis cs:0,0) -- (axis cs:0,1);
            \draw[dashed, gray!60] (axis cs:1,0) -- (axis cs:1,1);

            % Tracciamento del segnale h(t)
            \addplot[thick, nord10] coordinates {
                    (-1.5,0) (0,0) (0,1) (1,1) (1,0) (2.5,0)
                };
        \end{axis}
    \end{tikzpicture}
    \caption{Risposta impulsiva equivalente del sistema}
    \label{fig:Es_01_04_Esercizio1_RispostaImpulsivaEquivalente}
\end{figure}

Notiamo che possiamo quindi riscrivere $h(t)$ come segue:
\[
    h(t) = \text{rect} \left( t - \frac{1}{2} \right)
\]

E la \textbf{risposta in frequenza} risulta quindi essere:
\[
    \begin{aligned}
        H(f) &= \text{sinc}(f) e^{-j 2\pi f \frac{1}{2}} \\ 
             &= \text{sinc}(f) e^{-j\pi f}
    \end{aligned}
\]

\section{Esercizio 2}

Consideriamo ora un \textbf{sistema LTI} con la
seguente \textbf{risposta in frequenza}:
\[
    H(f) = j = e^{j \pi /2}
\]

Consideriamo inoltre un segnale di \textbf{ingresso sinusoidale}:
\[
    x(t) = A \sin (2\pi f_{0} t)
\]

\begin{exe}
    Vogliamo:
    \begin{enumerate}
        \item Determinare la \textbf{risposta del sistema} all'ingresso 
            dato sopra, $\forall A,f_{0}$.
        \item L'\textbf{energia} del segnale $y_{2}(t)$, uscita del sistema
            per l'ingresso $x_{2} = \frac{\sin 4\pi t}{2\pi t}$.
    \end{enumerate}
\end{exe}

Per quanto riguarda il \textbf{primo punto}, essendo il sistema \textbf{LTI},
sappiamo che l'uscita in risposta ad un segnale sinusoidale sarà
un \textbf{altro segnale sinosoidale} con \textbf{modulo} e 
\textbf{fase modificati}. Abbiamo quindi:
\[
    y(t) = A |H(f_{0})| \sin \left( 2\pi f_{0}t + \angle H(f_{0}) \right)
\]

Dobbiamo quindi calcolare $|H(f)|$ e $\angle H(f)$:
\[
    \begin{aligned}
        |H(f)| &= \left| e^{j \pi /2} \right| = 1 \text{ }\forall f \\ 
        \angle H(f) &= \frac{\pi}{2} \text{ }\forall f
    \end{aligned}
\]

E troviamo quindi:
\[
    \begin{aligned}
        y(t) &= A \sin (2\pi f_{0} t + \pi /2) \\ 
             &= A \cos (2 \pi f_{0}t)
    \end{aligned}
\]

Passiamo ora al secondo punto. 

\begin{remark}
    Possiamo sfruttare il \textbf{teorema di Parseval} e
    lavorare nel \textbf{dominio della frequenza} per
    \textit{evitare di calcolare esplicitamente} $y_{2}(t)$.
\end{remark}

Per cominciare, ricordiamo che
\[
    Y_{2}(f) = H(f) \cdot X_{2}(f)
\]
Calcoliamo quindi $X_{2}(f)$ osservando che $x_{2}(t) = 2 \text{sinc}(2t)$:
\[
    X_{2}(f) = 2 \cdot \frac{1}{2}\text{rect} \left( \frac{f}{2} \right) 
    = \text{rect} \left( \frac{f}{2} \right)
\]

Per cui:
\[
    Y_{2}(f) = j \cdot \text{rect} \left( \frac{f}{2} \right)
\]

A questo punto, possiamo applicare il \textbf{teorema di Parseval}
come anticipato sopra:
\[
    \begin{aligned}
        E_{y_{2}} &= \int_{-\infty}^{\infty} |y_{2}(t)|^2 dt \\
                  &= \int_{-\infty}^{\infty} y_{2}(t) \overline{y_{2}(t)} dt \\ 
                  &= \int_{-\infty}^{\infty} Y_{2}(f) \overline{Y_{2}(f)} df \\ 
                  &= \int_{-\infty}^{\infty} |Y_{2}(f)|^2 df
    \end{aligned}
\]

E troviamo infine:
\[
    \begin{aligned}
        E_{y_{2}} &= \int_{-\infty}^{\infty} |j \text{rect} \left( \frac{f}{2} \right) |^2 df \\ 
                  &= \int_{-\infty}^{\infty} |j|^2 \left| \text{rect} \left( \frac{f}{2} \right) \right|^2 \\ 
                  &= \int_{-\infty}^{\infty} 1 \cdot \text{rect} \left( \frac{f}{2} \right) df \\
                  &= \int_{-1}^{1} df \\ 
                  &= 2
    \end{aligned}
\]
